% Document class and packages
\documentclass[a4paper, 12pt]{report}
\usepackage{geometry}
\usepackage[utf8]{inputenc}
\usepackage{graphicx}
\usepackage{titlesec}
\usepackage{fontenc}
\usepackage{float}
\usepackage[swedish]{babel}
% For headers and footers
\usepackage{fancyhdr, lastpage}
\usepackage{fancyvrb}
% Math spesific packages
\usepackage{amsfonts}
\usepackage{amsmath}
\usepackage{amssymb}
% For reffereces displaying and coloring as well as for some document settings
\usepackage[]{hyperref}
% For math symbols and figures
\usepackage{tikz}
% For color boxes
\usepackage{xcolor}
\usepackage[most]{tcolorbox}
% For unit displaying
\usepackage{siunitx}
% For correct paragraph phrasing
\usepackage{parskip}

% For figure caption formatting
\usepackage[font=small,labelfont=bf,center]{caption}

% For correct fotnotes and citation structure
\usepackage[backend=bibtex8,style=verbose,citestyle=alphabetic,bibstyle=alphabetic]{biblatex}
\addbibresource{ref.bib}

% Background image of title page
\usepackage{eso-pic}
\newcommand\BackgroundIm{
\put(0,0){
\parbox[b][\paperheight]{\paperwidth}{
\vfill
\centering
\includegraphics[height=\paperheight,width=\paperwidth,
keepaspectratio]{background-1.png}
\vfill
}}}

% To fix footnote counter between files
\usepackage{chngcntr}
\counterwithout*{footnote}{chapter}

% Global options
\geometry{a4paper}
\graphicspath{{Images/}}
\hypersetup{
  linkcolor=black,
  citecolor={blue!50!black},
  filecolor={blue!50!black},
  urlcolor={blue!80!black},
  pdftitle={The AES encryption algorithm},
  pdfauthor={Gabriel Lindeblad},
  pdfcreator={Gabriel Lindeblad},
  pdfsubject={Advanced Encryption Standard},
  pdfstartpage=1,
  pdffitwindow=true,
  colorlinks=true,
  pdfpagemode=UseNone,
  }
\setcounter{tocdepth}{3}
\setcounter{secnumdepth}{3}
\pagestyle{fancy}
\rhead{\currentchapter}
\cfoot{\thepage}
\fancyheadoffset[L, RO]{0cm}


% New commands and changed commands
\newcommand{\currentchapter}{}
\newcommand{\mychapter}[1]{\chapter{#1}
\renewcommand{\currentchapter}{#1}}
\renewcommand{\bibname}{Källförteckning}
\renewcommand{\contentsname}{Inehållsförteckning}
\renewcommand{\listfigurename}{Figurer}
\newcommand{\fakechapter}[2]{%
  \par\setcounter{chapter}{#2}
  \chaptermark{#1}
  \rhead{ }
}
\newcommand*{\captionsource}[2]{%
  \caption[{#1}]{%
    #1%
    \\\hspace{\linewidth}%
    \textbf{Källa:} #2%
  }%
}

% Titel reformating
\titleformat{\chapter}[block]
  {\normalfont\huge\bfseries}{\thechapter}{0.5em}{\Huge}
\titlespacing*{\chapter}{0pt}{-25pt}{0pt}

\titleformat{\section}[block]
  {\normalfont\Large\bfseries}{\thesection}{0.5em}{\Large}
\titlespacing*{\chapter}{0pt}{-25pt}{0pt}

\titleformat{\subsection}[block]
  {\normalfont\large\bfseries}{\thesubsection}{0.5em}{\large}
\titlespacing*{\chapter}{0pt}{-25pt}{0pt}

\titleformat{\subsubsection}[block]
  {\normalfont\normalsize\bfseries}{\thesubsubsection}{0.5em}{\normalsize}
\titlespacing*{\chapter}{0pt}{-25pt}{0pt}


% Glossaries setup
\usepackage[nopostdot,nonumberlist,style=super,acronym,toc]{glossaries}
\setlength{\glsdescwidth}{0.8\hsize}
\makeglossaries

\renewcommand*{\glossaryname}{Begreppförklaring}
\renewcommand*{\acronymname}{Akronymer}
\renewcommand{\glsnamefont}[1]{\textbf{#1}}

%%%
\newglossaryentry{rsa}
{
    name={RSA},
    description={Rivest-Shamir-Adleman (RSA) är en av de mest välkända krypteringsalgoritmerna
    och var en av de första algorithmerna som byggde på en asymetrisk kryptering.
    RSA bygger på multiplikation av stora primatal där primtalen är
    nycklarna.\cite{rsa-ref}}
 }

\newglossaryentry{caesar}
{
    name={Caesarchiffer},
    description={Caesarchiffer är ett \gls{substitutionsskiffer},
    vilket helt enkelt bygger på att man byter ut varje bokstav i
    medelandet med en annan. Ersättnings bokstaven bestäms genom
    att man hoppar ett visst antal hopp i alphabetet som exempelvis
    3 hopp, vilket då innebär att ifall man har bokstaven a då skulle
    den bli ett d istället.\cite{caesar}}
 }

\newglossaryentry{xor}
{
    name={XOR},
    symbol={\(\oplus\)},
    description={XOR är en logisk operation inom datorvetenskap som fungerar ungefär som + uttrycket,
    med den enda skillnaden att {\(1 \oplus 1 = 0\)}. Detta samt att xor är en \gls{binär} operation, vilket
    innebär att termerna bara kan vara 0 eller 1 och resultatet det samma. Utöver XOR finns även
    \gls{or}, \gls{not} och \gls{and} bland annat.\cite{xor}} % Maybe needs to be refined
 }

\newglossaryentry{python}
{
    name={Python},
    description={Python är ett högnivå programmerings språk byggt på programmerings språket C.
    De är skapat av Guido van Rossum och släpptes i Februari 1991.\cite{python}}
 }

 \newglossaryentry{keystream}
 {
     name={Nyckelström},
     description={En nyckelström är i kryptografin en ström av \gls{pseudoslump}
     karaktärer som kan kombineras med exempelvis ett medelande för att producera
     en skiffertext.\cite{keystream}}
  }

  \newglossaryentry{pseudoslump}
  {
      name={Pseudoslump},
      description={Pseudoslump är en rad av nummer som kan se ut att vara helt
      slumpmässiga men har blivit framställda genom en upprepbar process.\cite{pseudoslump}}
   }

  \newglossaryentry{vscode}
  {
      name={VSCode},
      description={Visual Studio Code är en programutveklingsmiljö som är skapad av Microsoft.
      Det är ett öppet källkods projekt som är tillgängligt för det flesta operativsystem
      och kan användas för att skriva kod i flera olika språk.\cite{vscode-ref}}
   }

  \newglossaryentry{streamcipher}
  {
      name={Strömskiffer},
      description={Strömskiffer, ett symetriskt nyckel skiffer där man använder en
      \gls{pseudoslump}mässig skiffer ström (\gls{keystream}) som sedan en \gls{bit} i taget
      kombineras med de som ska krypteras. Den kombinerande operationen som används i
      strömskiffer är ofta en \gls{xor}-operation.\cite{streamcipher-ref}}
   }

  \newglossaryentry{bit}
  {
       name={Bit},
       description={En Bit är den minsta enheten av information som kan lagras i en dator.
       En bit kan endast ha två värden där den antingen är 0 eller 1, alltså ett \gls{binär}t
       värde. I datorvetenskap pratar man dock mer vanligen om ett \gls{byte} som är 8 bits.\cite{bit-ref}}
   }

  \newglossaryentry{byte}
  {
      name={Byte},
      description={En Byte består av 8 \gls{bit}s och är en enhet som används inom
      datorvetenskap. En byte kan ha 256 olika värden från 0 till 255, vilket är $2^8$ värden.
      Dessa värden representerar ofta tecken eller symboler som exempelvis bokstäver, siffror med mera.
      Tolkningen av vad en sekvens av bytes eller en enskild byte står för beror däremot på vilken
      teckenkodning som används. Exempel på teckenkodningar kan var \acrshort{ascii} och ISO-8859.\cite{byte-ref}}
   }

  \newglossaryentry{hashfunktion}
  {
      name={Hashfunktion},
      description={Hashfunktion är en funktion som man kan säga delar upp en viss datamängd
      och genom för en serie operationer som resulterar i hashtext av godkännd längd.
      Längden är samma för alla hastexter som använder samma funktion medans innehållet
      förändras så fort något över huvud taget ändras i datan som funktionen appliceras på.
      Användningsområdet för dessa funktioner är bland annat när man vill kunna verifiera
      medelanden eller information och försäkra sig om att ingen ändrat på medellandet efter att de skickats.
      Detta kan man då göra för att man vet att om man skör informationen genom samma hashfunktion
      borde resultatet vara identisk ifall informationen är oförändrad.\cite{hashfunktion-ref}}
   }

  \newglossaryentry{polyalphabetic-substitutionsskiffer}
  {
      name={Polyalphabetic substitutionsskiffer},
      description={polyalphabetic substitutionsskiffer bygger på att man använder flera olika \gls{substitutionsskiffer} för att
      på så sätt undvika en ut av de största svagheterna med \gls{substitutionsskiffer}. Detta då att dom lätt går att
      knäcka genom en \gls{frekvensanalys} då vissa bokstäver dyker upp mer frekvent i språket än andra. För att lösa
      detta så använder polyalphabetiska skiffer flera olika substitutionsskiffer som man byter mellan med en viss
      frekvens för att eliminera \gls{frekvensanalys}ens effektivitet.\cite{polyalphabetic-ref}}
   }

  \newglossaryentry{frekvensanalys}
  {
    name={Frekvensanalys},
    description={Frekvensanalys inom kryptografi är en metod för att knäcka ett \gls{substitutionsskiffer} genom att
    analysera frekvensen av bokstäver och utnyttja de faktum att en del bokstäver framkommer mer frekvent än andra
    i språket. På detta viset kan man då sedan lista ut vilka bokstäver som är vilka i det krypterade medellandet.\cite{frekvensanalys-ref}}
    }

  \newglossaryentry{substitutionsskiffer}
  {
      name={Substitutionsskiffer}, % fix this
      description={Ett Substitutionsskiffer är en typ av krypterings metod som bygger på att man
      byter ut delar av informationen man ska kryptera med exempelvis andra symboler med hjälp av en nyckel.
      Detta kan exempelvis vara bokstäver som byts ut mot andra bokstäver precis som i \gls{caesar} eller
      siffror som byts ut mot andra siffror.\cite{substitutionsskiffer-ref}}
    }

  \newglossaryentry{and}
  {
      name={AND},
      description={AND är en logisk operation inom datorvetenskap och matematik
      som tar två \gls{binär}a värden och ger till baka ett \gls{binär}t värd. Detta värde
      är 1 om och endast om båda värdena är 1 annars är värdet 0.\cite{logical-operators-ref}}
    }

  \newglossaryentry{or}
  {
      name={OR},
      description={OR är en logisk operation inom datorvetenskap och matematik
      som tar två \gls{binär}a värden och ger till baka ett \gls{binär}t värd. Detta värde
      är 1 om minst ett av värdena är 1 annars är värdet 0.\cite{logical-operators-ref}}
    }

  \newglossaryentry{not}
  {
      name={NOT},
      description={NOT är en logisk operation inom datorvetenskap och matematik
      som tar två \gls{binär}a värden och jämför dom. Det slutgiltiga värdet som ges tillbaka är
      1 om värdena inte är lika varandra och annars är värdet 0.\cite{logical-operators-ref}}
    }

  \newglossaryentry{binär}
  {
      name={Binär},
      description={Binär är ett begrepp som används inom datorvetenskap och matematik för att beskriva
      ett värde som endast kan vara två olika saker som exempelvis 0 eller 1, sant eller falskt. Detta innebär abstract
      binära tal bygger på talbasen 2, vilket skiljer sig från vad som ofta vanligen används i vardagen
      som är talbasen 10.\cite{binär-ref}}
    }

  \newglossaryentry{enigma}
  {
      name={Enigma},
      description={Enigma var en...
        .\cite{enigma-ref}}
    }

\newacronym{aes}{AES}{Advanced Encryption Standard}
\newacronym{ecb}{ECB}{Electronic Code Book läge}
\newacronym{cbc}{CBC}{Cipher Block Chaining läge}
\newacronym{ofb}{OFB}{Output Feedback läge}
\newacronym{des}{DES}{Data Encryption Standard}
\newacronym{iv}{IV}{Initialization Vector}
\newacronym{ssl}{SSL}{Secure Socket Layer}
\newacronym{tls}{TLS}{Transport Layer Security}
\newacronym{ascii}{ASCII}{American Standard Code for Information Interchange}
%%%

% Start document
\begin{document}

\newgeometry{top=0cm, bottom=2cm, left = 1.5cm, right = 1.5cm}

% Beginning title page
\begin{titlepage}
    % Inserting background images
    \AddToShipoutPicture*{\BackgroundIm}

    % School logo
    \begin{figure}[H]
    \includegraphics[width=0.3\textwidth]{skola_logga.png}
    \end{figure}

    % Aligning year to the right
    \raggedleft

    % Year
    \vspace{-3.5cm}
    {\large \textbf{2022/2023}}

    % Aligning to the left
    \raggedright

    % Title and sub title
    \vspace{3.5cm}
	{\huge\bfseries The AES encryption algorithm\par}
   ~{\large En undersökning av The \acrfull{aes}}

% Bottom of the page
    \vfill

  {Klass:} \\
    ~\textbf{NA20}

  {Handledare:} \\
    ~\textbf{Jimmy Nylén}

  {Författare:} \\
    ~\textbf{Gabriel Lindeblad}

  {Program:} \\
    ~\textbf{Naturvetenskapsprogrammet}

    \vspace{1cm}
	{\large \today\par}
    \vspace{-1cm}

% End of title page
\end{titlepage}

% Geometry of page borders for abstract
\newgeometry{top=5cm, bottom=2cm, left = 2cm, right = 2cm}

% Abstract
\chapter*{Abstract}
Lorem ipsum dolor sit amet, consectetur adipiscing elit. Ut lacinia ex eget sagittis congue.
Nullam cursus egestas dolor, suscipit gravida magna ultrices sit amet. Nullam placerat dui eu
arcu pharetra, sit amet tempor dolor convallis. Aenean sodales condimentum turpis, commodo maximus
augue. Aenean vel nibh dui. Pellentesque ex libero, lacinia nec mauris vel, convallis consectetur
felis. Maecenas ut nibh sed magna maximus imperdiet at id purus. In vel consequat metus. Donec non
tincidunt nunc. Sed pulvinar odio ut sapien vestibulum, quis mollis arcu tempor. Maecenas ut sem leo.
Sed leo risus, mollis eu ex vitae, feugiat consequat metus. Aenean interdum volutpat urna, nec
tempor mi accumsan quis. Morbi blandit maximus urna non aliquet aes.

% New page borders for bulk of the report
\newgeometry{top=2.5cm, bottom=2cm, left = 2cm, right = 2cm} % 2.25cm is the height of the header (is a checky solution that might need to bee removed if more headings are added)

% Inehållsförteckning
\tableofcontents

\clearpage

\printglossary

\printglossary[type=\acronymtype]

\mychapter{Inledning} % Mostly clear but might need some refinement
Kryptering, en bärande grundsten i dagens digitaliserade samhälle. Det är väggen mellan oss och resten
av världen, ett lås runt våra liv. Kryptering bygger på ett simpelt koncept, att dölja informationen
från all förutom den menade mottagaren. Ett koncept som har funnits med oss under stora delar av
människans historia och som än idag är en viktig del av vårt samhälle.\footfullcite{luciano1987cryptology}

Att dölja information har gått från de enkla metoder som användes redan för 2000 år sedan för att förmedla hemlig information med hjälp av
exempelvis \gls{caesar} till det komplexa och invecklade algoritmer som idag används för att skydda nästan all information som lagras och skickas
över internet.\footcite{luciano1987cryptology} Dessa komplexa algoritmer har utvecklats under en tid där datorer står som det dominerande informationshanteringsverktyget och
därför är det även en dator som används som huvudverktyg i denna rapport för att undersöka en av det vanligaste krypteringsalgoritmerna i dagens samhälle,
The \acrlong{aes}.

\section{Syfte} % Mostly clear but might need some refinement
Syftet med undersökning är att undersöka krypterings algoritmen \acrshort{aes},
för att utveckla en fördjupad förståelse för mer avancerade krypterings algoritmer.
Samt bygga en uppfattning om hur det på olika sätt går att implementera krypterings algoritmer
och vilken påverkan detta då får på deras prestanda och säkerhet.

\section{Frågeställningar} % Mostly clear but might need some refinement and reformulating
\begin{itemize}
    \setlength{\itemindent}{-1em}
    \item Hur påverkas tiden det tar att kryptera något mellan det olika nyckel längderna 128-bit,
          192-bit och 256-bit nyckel?

    \item Hur förändras krypterings tiden mellan de olika körlägena \acrshort{ecb}, \acrshort{cbc} \& \acrshort{ofb}?

    \item Hur påverkas skiffertexten av det olika körlägena \acrshort{ecb}, \acrshort{cbc} \& \acrshort{ofb}
          samt vilken betydelse får det ur ett säkerhetsperspektiv?
\end{itemize}

\section{Avgränsning} % Mostly clear but might need some refinement and reformulating

Denna undersökning är avgränsad till att endast undersöka \acrshort{aes} och dess
användning med fokus på hur nyckellängd och körläge påverkar krypteringstiden. Detta samt hur
den resulterande skiffer texten påverkas av det olika körlägena \acrshort{ecb}, \acrshort{cbc} \& \acrshort{ofb}
ur ett säkerhetsperspektiv.

Denna analys av algoritmens säkerhet tar inte hänsyn till faktorer så som möjliga attacker där
ibland exempelvis Brute-Force\footfullcite{kumar2011investigations}
\& Side-Channel\footfullcite{10.1007/978-3-642-04138-9_8} attacker. Undersökningen är även
begränsad till att endast utföras på en mjukvaruimplementering av \acrshort{aes}.

\chapter{Bakgrund}

\section{section 1}
abstract
\section{section 2}
abstract
\section{section 3}
abstract

\mychapter{Teori}


\section{Kryptering}


\section{Blockchiffer}


\subsection{Körlägen}


\subsubsection{ECB}
\acrfull{ecb} är en av det enklaste blockchiffer körlägena som finns.
\acrshort{ecb} i sig är ganska lätt att förstå och bygger i huvudsak bara på
att man delar upp den data man vill kryptera i delar kallade block och tar sedan varje
block för sig och kör genom algoritmen, vilket tydligt visas i
Figur \ref{fig:ecb-mode-enc} \& \ref{fig:ecb-mode-dec}.
\footfullcite{modesofoperation}

Figur \ref{fig:ecb-mode-enc} visar hur \acrshort{ecb} fungerar vid kryptering.
Här visas hur varje block för sig krypteras med hjälp av en blockchiffer algoritm
tillsammans med den givna nyckeln.

\begin{figure}[H]
    \includegraphics[width=\textwidth]{ECB_encryption.png}
    \captionsource{\acrlong{ecb} kryptering}{{\cite{ecb-mode-enc-ref}}}
    \label{fig:ecb-mode-enc}
\end{figure}

Figur \ref{fig:ecb-mode-dec} visar istället hur \acrshort{ecb} fungerar vid
dekryptering, vilken är en till stort sett identisk operation med det enda undantaget
att blockchiffret körs i dekrypterings läge istället för krypterings läge.

\begin{figure}[H]
    \includegraphics[width=\textwidth]{ECB_decryption.png}
    \captionsource{\acrlong{ecb} dekryptering}{{\cite{ecb-mode-dec-ref}}}
    \label{fig:ecb-mode-dec}
\end{figure}

På grund av \acrshort{ecb} körlägets simplicitet så finns det dock även ett ganska
stort problem med detta körläge. Det handlar om att \acrshort{ecb} inte på något
sätt förhindrar att två block med samma innehåll som krypteras inte resulterar i
ett identiskt krypterat block.\footcite{modesofoperation}

Vad detta innebär är att för större mängder data
är att det börjar bildas mönster i skiffertexten. Detta är något som väldigt
tydligt visar sig ifall man krypterar en bild, vilket går att se när man jämför
figur \ref{fig:pi-original} \& \ref{fig:pi-ecb}.
Det här faktumet är även varför \acrshort{ecb} inte är ett säkert körläge
och därför inte används näst intill aldrig i praktiken.\footcite{modesofoperation}

\subsubsection{CBC}
\acrlong{cbc} är ett av de mest vanligen använda körlägena för många blockchiffer.

\footcite{modesofoperation}

\begin{figure}[H]
    \includegraphics[width=\textwidth]{CBC_encryption.png}
    \captionsource{\acrlong{cbc} kryptering}{{\cite{cbc-mode-enc-ref}}}
    \label{fig:cbc-mode-enc}
\end{figure}

\begin{figure}[H]
    \includegraphics[width=\textwidth]{CBC_decryption.png}
    \captionsource{\acrlong{cbc} dekryptering}{{\cite{cbc-mode-dec-ref}}}
    \label{fig:cbc-mode-dec}
\end{figure}

\subsubsection{OFB}


\section{Symetrisk \& Asymmetrisk Kryptering}
Symetrisk och asymetrisk kryptering handlar om hur nycklar används i olika
krypteringsalgoritmer. För symetriska krypterings algoritmer så betyder detta
att samma nyckel är vad som används för både kryptering och dekryptering.

\footfullcite{symencrypt}

\section{AES}


\subsection{Finite Fields}


\subsection{AES S-Box}


\subsection{Struktur}


\subsubsection{SubBytes operation}
SubBytes operationen bygger på ...

\subsubsection{ShiftRows operation}


\subsubsection{MixColumns operation}


\subsubsection{AddRoundKey operation}


\subsection{Nyckel utökning}


\subsubsection{RotWord}


\subsubsection{SubWord}


\subsubsection{Rcon}


\subsection{AES-128bit}


\subsection{AES-192bit}


\subsection{AES-256bit}



\mychapter{Metod \& Genomförande} % (To be reviwed)
Metoden för denna undersökning bygger på en implementering av \acrshort{aes} i programmeringsspråket
\gls{python}. Detta tillsammans med ett antal konstruerade tester även dom implementerade i
\gls{python} är vad som använts för själva undersökningen av \acrshort{aes}. Koden
är skriven med hjälp av programmet \gls{vscode} och är byggd huvudsakligen för \gls{python} 3.10 men
på grund av att \gls{python} 3.11 släpptes innan undersökningen genomfördes så är de istället \gls{python} 3.11.0
som användes under undersöknings genomförandet.\footfullcite{python311}

\section{Implementering} % (To be reviwed)
Implementeringen av \acrshort{aes} är uppdelad i ett antal funktioner till stor del är baserat på
hur strukturen och uppdelningen av \acrshort{aes} beskrivs i \citetitle{daemen1999aes}\footcite{daemen1999aes}.
Några av det Huvudsakliga funktionerna av algoritmen är som följande:
\begin{itemize}
    \item \texttt{\nameref{sec:aes-subbytes}}
    \item \texttt{\nameref{sec:aes-shiftrows}}
    \item \texttt{\nameref{sec:aes-mixcolumns}}
\end{itemize}

Dessa funktioner används i varje runda och utgör grunden av algoritmen. Utöver detta finns
även \nameref{sec:aes-addroundkey} som används för att lägga till nyckeln i varje runda. AddRoundKey steget är dock inte
en funktion eftersom det enbart består ut av en \gls{xor}-operation. Utöver detta finns det även
en funktion som används för att expandera nyckeln som beskrivs i \nameref{sec:aes-key-expansion} samt
individuella funktioner för kryptering och dekryptering till de olika körlägen \nameref{sec:ecb}, \nameref{sec:cbc} och
\nameref{sec:ofb}.

Implementeringen använder sig av \gls{python} biblioteket NumPy för att hantera matriserna av \gls{byte}s och genomföra
matematiska operationer på dessa. Hela implementeringen går att hitta i appendix \ref{app:python} "\nameref{app:python}" samt
förklaras även i avsnittet \nameref{sec:aes}.

\section{Test Uppsättning} % (To be reviwed)
Test uppsättningen är uppdelad i det tre delarna \nameref{sec:nyckellangd-test}, \nameref{sec:körlages-test} och \nameref{sec:krypterings-test}. Två av delarna, \nameref{sec:nyckellangd-test} och
\nameref{sec:körlages-test} är skrivna som ett \gls{python} script vilket visas i appendix \ref{app:analyze}. När det gäller Krypteringstest så skiljer sig det från de två andra analyserna eftersom det
inte är en automatiserad process utan ett antal praktiska steg som genomfördes för att generera resultatet.

\subsection{Testmiljö} % Describe the test enviroment (might reqiuer extra glossary entrys for some of the terms)
Genomförandet av undersökningen skedde på en Windows dator med \gls{python} 3.11.0 installerat. I syfte att minimerar variabiliteten
mellan de olika omgångarna av tester då allt repeteras ett vist antal gånger så begränsades vissa delar av datorn och anpassningar gjordes i genomförandet.
De begränsningar som placerades på datorn bestod ut av att datorn CPU klockhastighet begränsades till sin grund hastighet för att ge
varje test ett likvärdigt utgångsläge som möjligt. Utöver detta så genomfördes även varje test separat utan parallellisering för att
eliminera risken av att testerna skulle kunna påverka varandra.

\subsection{Nyckellängds Test} % Explain the key length test
\label{sec:nyckellangd-test}
Nyckellängds testet bygger på att varje nyckellängd används för att kryptera en text fil på 1 MB. Detta genomförs 25 gånger för
varje nyckellängd och resultatet av varje test sparas sedan ...

\subsection{Körläges Test} % Explain the mode of operation test
\label{sec:körlages-test}


\subsection{Krypterings Test} % explain how the encryption test is setup
\label{sec:krypterings-test}
Krypterigns testet består utav att man krypterar en bild i bild formatet \gls{ppm} med hjälp av \acrshort{aes} i tre olika körlägen.
PPM formatet används då de tack vare sin enkla struktur kan genomgå en kryptering och fortfarande representeras som en bild
efteråt. Detta gör det då möjligt att visualisera resultatet av krypteringen. För att kryptera bilden används den implementerade
\acrshort{aes} algorithmen som körs i körlägena \nameref{sec:ecb}, \nameref{sec:cbc} och \nameref{sec:ofb} körlägen.

Före och efter krypteringen modifieras dock \gls{ppm} filen så att den går att öppna när den väl är krypterad. Detta görs genom att de fyra
första raderna i filen tas bort och sätts sedan tillbaka efter krypteringen. Detta gör att filen fortfarande kan öppnas som en bild
efter krypteringen.

\section{Genomförande} % Describe how the test was preformed (To be reviwed)
Först kördes Nyckellängdstestet och Körlägestestet genom att köra \gls{python} filen analyze.py. Där efter
genomfördes Krypterings testet. För krypteringstestet valdes först en lämplig bild ut för krypteringen. Sedan
konverterades bilden till fil formatet \gls{ppm} med hjälp av programmet \gls{gimp}. Där efter utförs följande kommandon i en terminal:

\begin{minted}{bash}
    $ head -n 4 pi.ppm > header.txt

    $ tail -n +5 pi.ppm > body.bin

    $ python3 encrypt.py 2b7e151628aed2a6abf7158809cf4f3c body.bin ECB

    $ cat header.txt body.bin.enc > pi-ecb.ppm
\end{minted}

Här används file header.txt som en temporär förvaringsplats för de första fyra raderna i \gls{ppm} filen. Detta
efter som de första fyra raderna innehåller informationen som berättar hur filen senare ska tolkas. Där efter används
tail kommandot för att ta bort de första fyra raderna i filen och skriva resten till en ny fil med namnet body.bin.
Filen body.bin krypteras sedan med hjälp av \gls{python} filen encrypt.py med nyckeln 2b7e151628aed2a6abf7158809cf4f3c.

Innehållet från den nya filen body.bin.enc sätt sedan ihop med innehållet av header.txt som infogas som de fyra första raderna samt filen döps om till pi-ecb.ppm.
Slutligen så konverteras även filen tillbaka till ett mer lätthanterligt format med hjälp av \gls{gimp}.
Processen beskriven ovan upprepades sedan för de två andra körlägen \nameref{sec:cbc} och \nameref{sec:ofb} med
den enda skillnaden att de istället för ECB används CBC och OFB som argument till encrypt.py filen (Det tredje terminal kommandot).

Därefter påbörjades \nameref{sec:nyckellangd-test} och \nameref{sec:körlages-test}, vilket utfördes genom att köra
\gls{python} filen analyze.py. Filen analyze.py genomförde då krypteringen av en textfil på 1MB fylld med slumpmässigt genererade \gls{byte}s med hjälp av \acrshort{aes}
implementeringen som går att se i appendix \ref{app:python}. Detta genomfördes i de tre olika körlägena \acrshort{ecb}, \acrshort{cbc} och \acrshort{ofb} samt de tre olika nyckellängderna 128bit, 192bit och 256bit.
För varje nyckellängd och körläge så genomfördes krypteringen 25 gånger. Resultatet exporterades därefter till 6 olika text filer, vilkas innehåll sedan sammanställdes i \nameref{chap:resultat}.

\mychapter{Resultat}
\label{chap:resultat}
\section{Nyckellängds Test}
\label{sec:nyckellangd}

\section{Körläges Test}
\label{sec:korlages}

\section{Krypterings Test}
\label{sec:krypterings-test}


\mychapter{Diskussion} % Diskussion av resultat, presentera och tolka resultatet samt för ett nyanserat resonemang om vad resultatet betyder. Diskutera även eventuella brister i experimentet och hur dessa kan förbättras.
\label{chap:discussion}

% Tolka resultatet från nyckellängds undersökningen
Utifrån resultatet från undersökningen kan man då se hur tiden det tar att kryptera en fil på 1MB ökar när man använder längre nycklar så som 192-\gls{bit} 256-\gls{bit} nycklar
i förhållande till en 128-\gls{bit} nyckel. Något som bland annat visar sig när man tittar på den procentuella tidskillnaden mellan de olika nyckellängderna. Där man kan se en
tids skillnad på 18,8\% mellan 128-\gls{bit} \& 192-\gls{bit} nyckeln och 27,8\% mellan 128-\gls{bit} \& 256-\gls{bit} nyckeln.
Anledningen till detta är något som högst troligen beror på att antalet krypteringsrundor som genomförs för det olika nyckellängderna ökar när nyckeln blir längre.
Något som då höjer antalet operationer som genomförs för varje 16-\gls{byte} block som krypteras, vilket i sin tur höjer den totala krypterings tiden.

% Tolka resultatet från körläges undersökningen
Resultatet visar även att skillnaden i krypteringstid mellan olika körlägen är väldigt liten. Mellan \acrshort{ecb} \& \acrshort{cbc} var skillnaden 0,3\% och mellan \acrshort{ecb}
\& \acrshort{ofb} var skillnaden 0,4\%. Något som inte riktigt var väntat då det för både \acrshort{cbc} och \acrshort{ofb}
genomförs en ytterligare operation mellan varje 16-\gls{byte} block som krypteras för att länka ihop de olika blocken. Men samtidigt så skulle anledning till att skillnaderna är
så små kunna vara att \gls{xor}-operationen som används för att länka ihop blocken är en operation som är relativt lätt och snabb för en dator att genomföra, samtidigt som
operationen även bara behöver genomföras en gång för varje 16-\gls{byte} block som krypteras.

I resultatet av körläges testet kan man dock även se en viss spridning i tiden som det tar att kryptera en fil mellan det olika omgångarna.
Något som visar på en viss osäkerhet då \acrshort{ecb} har en högsta tid som är högre än både den högsta tiden för \acrshort{cbc} och \acrshort{ofb}. Vilket då skulle kunna antyda
att de skillnader som mätts upp mellan medelvärdena kan vara något som är inom felmarginalen för mätningen. Detta innebär då att det inte går att säga med säkerhet att
det finns någon större skillnad i krypterings tiden mellan det olika körlägena utifrån denna undersökning.

% Tolka resultatet från krypteringsmetod undersökningen
Tittar man sedan istället på hur säkerheten påverkas av det olika körlägena så kan man ganska tydligt se att \acrshort{ecb} är det mest osäkra körläget för större data mängder.
Detta då man som visas i resultatet kan se hur trots kryptering man fortfarande kan se spår av den ursprungliga bilden i informationen som krypterats. Något som inte
går att göra när bilden istället krypteras med hjälp av \acrshort{cbc} eller \acrshort{ofb}.

% Utvärdera metoden (nyanserat resonemang)



\section{Felkällor} % Variabilitet i klockhastighet, Ogämnheter i implementeringen som resulterar i missvisande tider, osäkerhet host mätmetoden, fel i implementeringen,


\section{Förbättringar} % Öka antalet reppetitioner för att minskar risken för fel, använda andra metoder för att mäta tiden, testa olika nycklar för att se hur de påverkar,


\section{Slutsats}



\fakechapter{Källförteckning}{7}
\addcontentsline{toc}{chapter}{Källförteckning}
\printbibliography[title=Källförteckning]

\clearpage

\fakechapter{Figurer}{8}
\addcontentsline{toc}{chapter}{Figurer}
\listoffigures

\begin{figure}
  \centering
  \includegraphics[width=0.6\textwidth]{pi.png}
  \caption{Orginal bild}
  \label{fig:pi-original}
\end{figure}

\begin{figure}
  \centering
  \includegraphics[width=0.6\textwidth]{pi-ecb.png}
  \caption{Efter ECB Kryptering}
  \label{fig:pi-ecb}
\end{figure}

\begin{figure}
  \centering
  \includegraphics[width=0.6\textwidth]{pi-cbc.png}
  \caption{Efter CBC Kryptering}
  \label{fig:pi-cbc}
\end{figure}

\begin{figure}
  \centering
  \includegraphics[width=0.6\textwidth]{pi-OFB.png}
  \caption{Efter OFB Kryptering}
  \label{fig:pi-ofb}
\end{figure}

\begin{figure}
  \centering
  \includegraphics[width=0.6\textwidth]{AES_(Rijndael)_Round_Function.png}
  \captionsource{Uppställning av vanliga rundor}{{\fullcite{rounds-picture}}}
  \label{fig:round-function}
\end{figure}

% End document
\end{document}